
\section{Calibration de l'axe spectral (\texorpdfstring{$\lambda \rightarrow \mathrm{cm}^{-1}$}{lambda->cm-1}) (Section 4)}

\subsection{Objectif}
Assurer une conversion précise \emph{pixel $\rightarrow$ longueur d'onde $\rightarrow$ nombre d'onde (cm$^{-1}$)} et une estimation de la résolution spectrale. La calibration fournit un mapping stable, des incertitudes, et un contrôle dérive.

\subsection{Principes physiques}
\textbf{Grating} (réseau) : $m\lambda = d(\sin \alpha + \sin \beta)$, avec ordre $m$, pas $d$, angles d'incidence $\alpha$ et de diffraction $\beta$. Autour d'un pixel central $p_0$ (géométrie quasi-Littrow), on linéarise la dispersion et on retient un modèle polynômial:
\begin{equation}
\lambda(p) = a_0 + a_1 (p-p_0) + a_2 (p-p_0)^2 + a_3 (p-p_0)^3.
\end{equation}
\textbf{Conversion Raman} : pour une excitation à $\lambda_L$,
\begin{equation}
\nu~(\mathrm{cm}^{-1}) = \left(\frac{1}{\lambda_L(\mathrm{nm})} - \frac{1}{\lambda_S(\mathrm{nm})}\right) \times 10^7.
\end{equation}

\subsection{Incertitudes (propagation)}
Soit $\sigma_{\lambda_S}$ l'incertitude sur $\lambda_S$ issue du fit et $\sigma_{\lambda_L}$ celle de la ligne laser:
\begin{equation}
\sigma_{\nu}^2 \simeq \left(\frac{10^7}{\lambda_S^2}\right)^2 \sigma_{\lambda_S}^2 + \left(\frac{10^7}{\lambda_L^2}\right)^2 \sigma_{\lambda_L}^2.
\end{equation}
On trace $\sigma_{\nu}$ pour documenter la précision effective en cm$^{-1}$.

\subsection{Standards et cibles}
\begin{itemize}[nosep]
\item \textbf{Rayleigh} ($\nu=0$) pour ancrer le zéro (si observable/notch mesurable).
\item \textbf{Gaz néon/argon} : raies étroites pour dispersion $\lambda(p)$.
\item \textbf{Solide de référence} (Si $\sim$520.7~cm$^{-1}$, diamant 1332~cm$^{-1}$) pour vérifier la conversion Raman et la résolution.
\end{itemize}

\subsection{Algorithme de calibration}
\begin{enumerate}[nosep]
\item Acquérir $K$ spectres de référence (temps court, plusieurs accumulations).
\item Détecter les raies et \emph{centroïder} par fit Voigt ($\lambda_{S,k} \pm \sigma_k$).
\item Résoudre le mapping $\lambda(p)$ par régression pondérée (poids $1/\sigma_k^2$); choisir l'ordre $1$--$3$ par AIC/BIC.
\item Convertir en $\nu$ via $\lambda_L$; vérifier que les raies connues retombent à $\nu$ attendus (RMSE).
\item Estimer la \textbf{résolution} $\delta \nu$ par largeur (FWHM) de raies étroites (déconvolution instrumentale si besoin).
\item Exporter \texttt{calibration.json}: coefficients $(a_0..a_3)$, $p_0$, $\lambda_L$, RMSE, $\delta\nu$, validité (plage de $p$).
\end{enumerate}

\subsection{Contrôle-qualité et dérive}
\begin{itemize}[nosep]
\item \textbf{Vérif. journalière}: mesurer 1--2 raies, recalculer l'offset; alerter si $\Delta \nu > \tau_{\mathrm{drift}}$ (ex.: $>0.5$~cm$^{-1}$).
\item \textbf{Température}: consigner $T$ et corriger si une dépendance $\Delta \nu (T)$ est observée.
\item \textbf{Notch/Edge}: si filtre déplace/élargit Rayleigh, ignorer la région masquée et ancrer sur raies proches.
\end{itemize}

\subsection{Pseudo-code}
\begin{lstlisting}[basicstyle=\ttfamily\small]
def calibrate_axis(calib_spectra, lambda_L_nm, p0):
    # 1) peak picking + Voigt centroid
    peaks = []
    for spec in calib_spectra:
        for pk in detect_lines(spec):
            lam, sig = fit_voigt_centroid(spec, pk)
            peaks.append({"p": pk.pixel, "lambda": lam, "sigma": sig})
    # 2) weighted polynomial fit lambda(p)
    X = design_matrix([r["p"]-p0 for r in peaks], order=3)
    y = np.array([r["lambda"] for r in peaks])
    w = 1/np.maximum(np.array([r["sigma"] for r in peaks])**2, 1e-18)
    a = weighted_ls(X, y, w, order_select="aic")
    # 3) resolution from narrow lines (FWHM)
    delta_nu = estimate_resolution_cm1(calib_spectra, lambda_L_nm, a, p0)
    # 4) RMSE check
    rmse = compute_rmse(peaks, a, p0)
    return {"a":a, "p0":p0, "lambda_L":lambda_L_nm, "delta_nu":delta_nu, "rmse":rmse}
\end{lstlisting}

\subsection{Tests \& DoD}
\begin{itemize}[nosep]
\item RMSE $\le$ 0.2--0.5~cm$^{-1}$ sur la plage utile; $\delta \nu$ cohérent avec la grille/fente.
\item Vérification croisée sur Si/diamant: écart $<0.5$~cm$^{-1}$.
\item Export \texttt{calibration.json} versionné et relu au runtime; \texttt{spectral\_resolution\_cm1} renseigné en méta.
\end{itemize}
